\documentclass[11pt,a4paper]{article}
\usepackage[top=1.00in, bottom=1.0in, left=1in, right=1in]{geometry}
\usepackage{graphicx}
\usepackage{sectsty,setspace,natbib,wasysym} 
\usepackage{natbib}
\usepackage{parskip}

\begin{document}
\bibliographystyle{/Users/Lizzie/Documents/EndnoteRelated/Bibtex/styles/naturemag}
\hspace{-5ex} \includegraphics[width=0.4\textwidth]{/Users/Lizzie/Documents/Professional/images/letterhead/ubc/Faculty of forestry 2026short.png}
\pagenumbering{gobble}

\vspace{2ex}
\begin{small}
\noindent 2424 Main Mall \\
\noindent Vancouver, BC Canada V6T 1Z4\\
\noindent Ph: 604.827.5246\\
\end{small}
\\
\pagenumbering{gobble}
\noindent Dear Dr. Polle:

Please consider our revised manuscript (NPH-V-2026-58419), entitled `The problem and promise of modeling plant leafout under warmer winters,' as a Viewpoint for  \emph{New Phytologist}. % The piece focuses on chilling---the concept of cool temperatures required to break dormancy---working molecular biology forest tree and crop modeling, and research focused on forecasting with climate change. %---to uncover the problems that may limit the potential for rapid progress through new experimental work. % As such we leverage existing data from different fields to show how common some of the unspoken problems with current approaches are.

This paper addresses why current models of leafout timing with climate change often make wildly different predictions for the same system and temperature regime. We argue that progress has stalled because we do not understand how winter temperatures affect leafout. Models today estimate the effect of winter temperatures via an unobserved process called chilling. This has produced fundamental problems that the field has struggled to admit or address. In this Viewpoint, we bring a cross-disciplinary and solutions-focused perspective on chilling, bridging across plant physiology, molecular biology, forest tree and crop modelers to highlight the problems that may limit the potential for rapid progress, and show how to fix them.

Comments from two referees have greatly improved the manuscript, leading us to add new text and figures, as well as make edits throughout for clarity. We have included more detail on new findings from the molecular and cellular literature, fully revising the related figure and providing complete references and detailed information for this figure in the Supporting Information. We have also added more information on models of chilling and a related new figure (Figure 2). Figure 3  has also been completely overhauled based on feedback, and we made edits to improve both other figures. We believe our revised manuscript is much improved and provides a motivating path forward for the study of chilling, and we detail these changes in our point-by-point response letter. 

\vspace{1ex}\\
Sincerely,\\

\includegraphics[scale=1]{/Users/Lizzie/Documents/Professional/Vitas/Signatures/SignatureLizzieSm.png} \\

\noindent Elizabeth M Wolkovich\\
Professor of Forest \& Conservation Sciences\\ 
\end{document} 